% page 344 of the facsimile (image p-343 of the scan)
% block 152.4 x 242.0 mm ; left margin 8.0 mm
\pgc{-12.700mm}{0.000mm}{
\l{\cel{3}336}
\l{}
\l{\sou{lex}iko. Vortolibro qua indikas lakonike (generale : perun}
\l{\cel{3}vorto) la senco di la vorti di altra linguo. - DEFIRS.}
\l{}
\l{\sou{lexikografo.}\cel{2}La persono qua exploras,studias, la origino}
\l{\cel{3}e la senci di la vorti di linguo. - DEFIS.}
\l{}
\l{\sou{lexikografio}. Ensemblo ek la laboruro de lexikografo.}
\l{}
\l{\sou{lexikologio.}\cel{2}Cienco pri la origino e la senci di la vorti}
\l{\cel{3}di linguo.}
\l{}
\l{\sou{lezar}. (\sou{trans}.)I. (\sou{fiziol}.)Igar ne-uzebla parte. - II.}
\l{\cel{3}(Yuro-cienco). Detrimentar per desegardo di (la yuri).}
\l{\cel{3}- DEFIS.}
\l{}
\l{\sou{li.}\cel{2}(\sou{gram}.)Pronomo personala, triesma persono, pluralo,}
\l{\cel{3}formo epicena. -}
\l{}
\l{\sou{liano}.\cel{2}\sou{(bot.})\cel{2}Planto sarmentoza, klimera, di la foresti}
\l{\cel{3}di Amerika, qua plekte kaptas la arbori e formacas den-}
\l{\cel{3}saji nedesintrikebla. - DEFIRS.}
\l{}
\l{\sou{liaso}. (\sou{geol.)}\cel{2}Tereno kalkoza, marnoza, arjiloza. - DEFIS.}
\l{}
\l{libacionar. (netrans.) (religio antiqua). Disvarsar adsule,}
\l{\cel{3}honorize di deajo, mielo, lakto, oleo e precipue vino}
\l{\cel{3}pura. - DEFIS.}
\l{}
\l{\sou{libelo}. (\sou{tekn.}) Instrumento uzata por trasar lineo, plano,}
\l{\cel{3}horizontala, o por mezurar, relate plano horizontala,}
\l{\cel{3}la elevaciono di lineo, di plano qua esas paralela a lu.}
\l{}
\l{\sou{libelulo}. (\sou{zool.}) Insekto nevroptera, kun korpo dina ed}
\l{\cel{3}ali diafana. - L.\cel{1}\sou{libellula.}\cel{2}- DEFIS.}
\l{}
\l{\sou{libera}. I. (\sou{aludante individuo)}. Qua ne apartenas a mastro.}
\l{\cel{3}- II. (\sou{aludante statano)}. Qua ne dependas de autoritato}
\l{\cel{3}arbitriala. - III. (\sou{aludante stato}).Qua ne esas sklav-}
\l{\cel{3}igita da lando stranjera. - IV. (\sou{aludante psiko)}. Di qua}
\l{\cel{3}la volo povas determinar su sen subisar koakto. -}
\l{\cel{3}V. (\sou{aludante la korpo, la spirito, la kordio di la homo)}.}
\l{\cel{3}Di qua la agado ne renkontras obstaklo. - VI. (\sou{aludante}}
\l{\cel{3}la\cel{1}\sou{kozi)}. Qua ne prizentas obstaklo.}
\l{}
\l{\sou{libertina}. I. (\sou{olim).}\cel{2}I. Qua su liberigis de la autoritato}
\l{\cel{3}di la religio kristana, di la kredi, di la diciplino.}
\l{\cel{3}- II. Qua su liberigis de omna regulo, de omna autori-}
\l{\cel{3}tato. - III. Di qua la mori sexuala esas sen-regulo.}
\l{\cel{3}- DEFIS.}
\l{}
\l{\sou{libro.}\cel{2}I. Reprodukturo imprimura di la verko da autoro,}
\l{\cel{3}ek paperfolii asemblita aden kayeri quin onu inter-}
\l{\cel{3}asemblas. - II. La verko tale reproduktita. - III. Sin-}
\l{\cel{3}gla ek la parti di ula verki. - EFIS.}
}
